% Architecture guide for ../http-client.cpp.
% Build: pdflatex -interaction=nonstopmode -halt-on-error http-client.tex
\documentclass[11pt,a4paper]{article}
\usepackage[T1]{fontenc}
\usepackage[utf8]{inputenc}
\usepackage{lmodern}
\usepackage[margin=25mm]{geometry}
\usepackage{microtype}
\usepackage{booktabs}
\usepackage{listings}
\usepackage{xcolor}
\usepackage[hidelinks]{hyperref}
\setlength{\parindent}{0pt}
\setlength{\parskip}{6pt}
\lstset{basicstyle=\small\ttfamily,columns=fullflexible,keepspaces=true,
  breaklines=true,frame=single,showstringspaces=false,
  commentstyle=\color{gray},keywordstyle=\color{blue!60!black}}
\newcommand{\code}[1]{\texttt{\detokenize{#1}}}
\title{Architecture of the Concurrent UM-Bridge C++ Client}
\author{\texttt{umbridge\_mox/sandbox/C++/hpc/http-client.cpp}}
\date{}
\begin{document}
\maketitle

\section{Purpose and execution model}

The client evaluates several independent inputs through a UM-Bridge load
balancer. Each job owns a separate \code{umbridge::HTTPModel}, a separate
HTTP connection context, and a randomly generated session identifier. This
lets the balancer associate each job with its own backend model server.
The client uses the implementation in \code{umbridge_mox/umbridge/lib}.

The main thread follows three phases:
\begin{enumerate}
  \item Create all asynchronous workers, assigning an input and a job index
        to each one.
  \item Consume every worker's future and store successful results in
        \code{outputs[i]}.
  \item After every worker has finished, print each successful job number
        and its corresponding output, in increasing job order.
\end{enumerate}
Progress diagnostics such as ``Creating client'' may appear during execution.
Numerical results are printed only in the final phase. Errors are reported
separately to standard error.

\begin{center}
\fbox{\begin{minipage}{0.90\linewidth}
\centering
Main thread: validate arguments and prepare inputs\\[3pt]
$\downarrow$ launch all workers with \code{std::launch::async}\\[3pt]
Worker 0 \quad Worker 1 \quad $\cdots$ \quad Worker $N-1$\\
Each worker: create session $\to$ validate $\to$ evaluate $\to$ close\\[3pt]
$\downarrow$ return outputs or propagate exceptions through futures\\[3pt]
Main thread: collect all futures $\to$ print stored results
\end{minipage}}
\end{center}

\section{Data model and responsibilities}

\code{ModelInput} and \code{ModelOutput} both alias
\code{std::vector<std::vector<double>>}. The outer vector represents blocks;
each inner vector holds the scalar values in one block. These names clarify
which side of an evaluation a value belongs to without introducing different
C++ types.

\begin{center}
\small
\begin{tabular}{@{}p{0.32\linewidth}p{0.60\linewidth}@{}}
\toprule
Component & Responsibility \\
\midrule
\code{main} & Validate arguments, prepare inputs, launch workers, collect
results and print them after completion. \\
\code{evaluate_job} & Own one session and carry out one validated evaluation. \\
\code{SessionCloser} & Attempt session closure when its scope ends, including
when an exception unwinds the worker. \\
\code{create_session_id} & Generate a 32-character hexadecimal session ID from
128 random bits. \\
\code{validate_address} & Check server-address syntax before making requests. \\
\code{validate_blocks} & Compare input or output block dimensions with model
metadata. \\
\code{to_string} & Format one numeric block for printing. \\
\code{print_help} & Describe the command line and execution behavior. \\
\bottomrule
\end{tabular}
\end{center}

\section{Launching independent workers}

For job index $i$, the current example prepares a single input block
$[100+i,18]$ and launches the worker as follows:
\begin{lstlisting}[language=C++]
ModelInput inputs = {{100.0 + static_cast<double>(i), 18.0}};
jobs.emplace_back(std::async(std::launch::async, evaluate_job,
                             host, model_name, inputs, i));
\end{lstlisting}
The explicit \code{std::launch::async} policy starts an asynchronous task.
Without this policy, \code{std::async} is allowed to defer execution until a
future is consumed, which would undermine the intended concurrency.

Task arguments are copied into the asynchronous operation. In particular,
\code{inputs} is passed by value: destroying the local input variable at the
end of a loop iteration does not invalidate a worker's data. The host and
model-name arguments are also retained by the asynchronous operation for the
call's lifetime.

The vector \code{jobs} contains \code{std::future<ModelOutput>} objects. It
does not contain models or cleanup guards. Consequently, vector reallocation
cannot invalidate references between an HTTP model and its guard.
Each worker constructs these objects on its own stack.

The example creates one asynchronous task per requested job; it does not use a
bounded client worker pool. For very large job counts, thread and connection
resources must be considered separately from the limit on backend PBS jobs.

\section{One worker and its session lifetime}

\code{evaluate_job(host, model_name, inputs, i)} executes the following steps:
\begin{enumerate}
  \item Generate a session ID and construct an \code{HTTPModel}. Construction
        discovers the available models and obtains the selected model's
        capabilities. Against the load balancer, model-specific discovery may
        wait for an available scheduler slot and a ready backend server.
  \item Construct a \code{SessionCloser} referring to that model.
  \item Verify support for evaluation and request the input and output sizes.
        All these operations use the same session and therefore reuse the same
        backend job.
  \item Validate the caller-supplied input dimensions, evaluate with the default
        model configuration, and validate the returned output dimensions.
  \item Explicitly call \code{CloseSession()} and return the output blocks.
\end{enumerate}
The job index is used for progress diagnostics; the input data is supplied by
\code{main}, rather than being generated inside the worker.

\subsection{Why declaration order matters}
\begin{lstlisting}[language=C++]
umbridge::HTTPModel client(host, model_name, false,
                            httplib::Headers(), create_session_id());
SessionCloser session_closer(client);
\end{lstlisting}
C++ destroys local objects in reverse declaration order. The guard is therefore
destroyed before \code{client}, so its borrowed reference remains valid during
cleanup. The guard is noncopyable, preventing an accidental second cleanup
owner. It does not own the model itself.


On success, explicit closure happens before the worker returns. In the local
UM-Bridge implementation, successful closure clears the session ID, making the
guard's subsequent call a no-op. If an operation throws after the guard has
been created, its destructor attempts closure during stack unwinding. It
reports standard cleanup exceptions rather than replacing the original
exception. If model construction itself fails, no guard exists yet; the local
\code{HTTPModel} constructor performs its own best-effort session cleanup.

Explicit closure is part of successful worker completion. If closure throws,
the worker reports failure through its future even if evaluation produced a
numerical result; the guard attempts closure again during unwinding.
Forced process termination does not run scope guards, and a failed network
cleanup request cannot guarantee that a remote job has ended.

\subsection{Why the cleanup guard has no defaulted move operations}
\code{SessionCloser} is the only class defined in \code{http-client.cpp}.
Its user-declared destructor and deleted copy operations suppress implicit
move generation. Adding a defaulted move constructor would be unsafe: moving
the reference member would leave both guards referring to the same model.
Destruction of the source guard could then close the session before the
destination guard's scope ends. A defaulted move-assignment operator would be
deleted by the compiler because reference members cannot be rebound.

The guard therefore remains nonmovable. Supporting moves would require a
different representation and custom operations that transfer cleanup
responsibility and disarm the source. The current worker-local lifetime does
not require such transfers. \code{ModelInput} and \code{ModelOutput} are aliases
for standard vectors, whose move operations are already provided by the
standard library; they introduce no additional classes requiring declarations.

\section{Results, ordering, and exceptions}

Each asynchronous worker supplies either a \code{ModelOutput} or an exception
through its future. The main thread consumes futures in job order:
\begin{lstlisting}[language=C++]
std::vector<ModelOutput> outputs(num_jobs);
for (int i = 0; i < num_jobs; ++i) {
  try {
    outputs[i] = jobs[i].get();
  } catch (const std::exception &error) {
    // Report this job and continue observing the remaining futures.
    failed = true;
  }
}
// Print outputs only after the collection loop has ended.
\end{lstlisting}
Calling \code{get()} waits if necessary, transfers the result, or rethrows the
worker's exception. Each future is consumed once. Waiting first for job zero
does not serialize evaluations: all workers were launched earlier and continue
running independently. Each worker closes its session before returning,
regardless of when \code{main} reaches its future. This is essential for
releasing scheduler slots while other jobs are still waiting.

Only the main thread writes to or reads the aggregate \code{outputs} vector.
No mutex is needed for this vector. A failed job leaves its entry empty;
validated successful jobs have at least one output block. The printing loop
skips failed entries, prints the job index, and formats every block. Multiple
blocks are separated by \code{ | }; values within a block are comma-separated.
For example:
\begin{lstlisting}
Job 0 output: 254.000000
Job 1 output: 256.000000
Job 2 output: 258.000000
\end{lstlisting}
A job failure does not cancel other jobs. The main thread observes all futures,
prints successful outputs after collection, and returns failure if any job
failed. Worker progress and error messages use \code{std::osyncstream} to keep
concurrent messages from interleaving at the character level.

\section{Interaction with the load balancer and PBS}

Client concurrency and scheduler capacity are separate. The client can launch
50 workers while the balancer is configured with
\code{--max-pbs-jobs=20}. The balancer admits at most 20 outstanding PBS jobs
owned by that balancer, including queued jobs. Additional session requests
wait for a slot. A worker releases its session as soon as its work is complete,
allowing a waiting worker to proceed. Actual running capacity also depends on
the cluster scheduler, or on \code{MAX_JOBS} in the fake PBS implementation.

The balancer's \code{SchedulerTaskQueue} dispatches HTTP connections
independently. A fixed HTTP pool could otherwise fill with new-session
requests waiting for slots, blocking the evaluation or session-close requests
needed to release those slots. The PBS limit bounds backend jobs; it does not
bound frontend HTTP threads. The client itself submits no \code{qsub} or
\code{qdel} commands: scheduler ownership belongs to the balancer.

\section{Building, running, and adapting the example}

Build from \code{umbridge_mox/sandbox/C++/hpc}:
\begin{lstlisting}[language=bash]
make http-client LoadBalancer
\end{lstlisting}
The build uses C++20 for \code{std::osyncstream}, thread support, and OpenSSL
for HTTPS-enabled HTTP clients.

From \code{umbridge_mox/sandbox/C++/test}, use separate terminals:
\begin{lstlisting}[language=bash]
# Terminal 1: use the bundled fake PBS commands.
PATH="../hpc/bin:$PATH" ./LoadBalancer --scheduler=pbs --max-pbs-jobs=20

# Terminal 2: request 50 independent evaluations.
./http-client http://localhost:4242 50
\end{lstlisting}
The job-count argument is optional and defaults to five. Job indices start at
zero. Use \code{./http-client --help} for syntax. Real-cluster deployments use
the site's PBS commands and adapted job scripts instead of the fake commands.

To provide application-specific inputs, replace the input construction in the
launch loop with a lookup in a dataset or another computation. Keep each
worker's model and session independent. To store results, add persistence after
the result-collection loop, using \code{outputs[i]} as the result belonging to
input $i$ and accounting for failed entries. This preserves parallel execution
and the association between inputs and outputs.

\end{document}
